\setcounter{section}{10}






  \zwischenueberschrift{Noethersche Moduln}

Wir wollen zeigen, dass für einen noetherschen Ring $R$ und einen endlich erzeugten $R$-Modul jeder $R$-Untermodul wieder endlich erzeugt ist. Solche Moduln nennt man noethersch.


\inputdefinition
{}
{





Es sei $R$ ein
\definitionsverweis {kommutativer Ring}{}{}
und $M$ ein
$R$-\definitionsverweis {Modul}{}{.}
Dann heißt $M$ \definitionswort {noethersch}{,} wenn jeder
$R$-\definitionsverweis {Untermodul}{}{}
von $M$
\definitionsverweis {endlich erzeugt}{}{}
ist.





}

Für

\mavergleichskette
{\vergleichskette
{ M
}
{ = }{ R
}
{  }{
}
{    }{
}
{   }{
}
}
{}{}{}
stimmt dies mit der Definition eines noetherschen Ringes überein, da ja die $R$-Untermoduln von $R$ gerade die Ideale sind.

In den folgenden Aussagen verwenden wir folgende Sprech- bzw. Schreibweise.


\inputdefinition
{}
{





Es sei $R$ ein
\definitionsverweis {kommutativer Ring}{}{}
und seien
\mathl{M_1,M_2,M_3}{}
$R$-\definitionsverweis {Moduln}{}{.}
Man nennt ein Diagramm der Form

\mathdisp {0 \, \stackrel{  }{\longrightarrow} \, M_1   \, \stackrel{  }{ \longrightarrow}   \,  M_2   \, \stackrel{  }{\longrightarrow} \,  M_3  \,\stackrel{  } {\longrightarrow}  \, 0} {  }

eine \definitionswort {kurze exakte Sequenz}{} von $R$-Moduln, wenn $M_1$ ein $R$-Untermodul von $M_2$ ist, und wenn $M_3$ ein Restklassenmodul von $M_2$ ist, der isomorph zu
\mathl{M_2/M_1}{} ist.





}

Die Exaktheit bedeutet, dass an jeder Stelle die Beziehung

\mavergleichskettedisp
{\vergleichskette
{  \operatorname{kern} \varphi_{i+1}
}
{ =} {  \operatorname{bild} \varphi_i
}
{ } {
}
{ } {
}
{ } {
}
}
{}{}{}
gilt, wenn $\varphi_i$ die
$R$-\definitionsverweis {Modulhomomorphismen}{}{}
bezeichnet.





\inputfaktbeweis
{Modultheorie (kommutative Algebra)/Noethersche Moduln/Kurze exakte Sequenz/Äquivalentes Kriterium/Fakt}
{Lemma}
{}
{



\faktsituation {}
\faktvoraussetzung {Es sei $R$ ein kommutativer Ring und

\mathdisp {0 \longrightarrow M_1 \longrightarrow M \longrightarrow M_3 \longrightarrow 0} {  }

eine
\definitionsverweis {kurze exakte Sequenz}{}{}
von $R$-Moduln.}
\faktfolgerung {Dann ist $M$ genau dann
\definitionsverweis {noethersch}{}{,}
wenn sowohl $M_1$ als auch $M_3$ noethersch sind.}
\faktzusatz {}
\faktzusatz {}





}
{





Es sei zunächst $M$ noethersch, und

\mavergleichskette
{\vergleichskette
{ U
}
{ \subseteq }{ M_1
}
{  }{
}
{    }{
}
{   }{
}
}
{}{}{}
ein Untermodul. Dann ist $U$ direkt auch ein Untermodul von $M$, also nach Voraussetzung endlich erzeugt. Es sei nun

\mavergleichskette
{\vergleichskette
{ V
}
{ \subseteq }{ M_3
}
{  }{
}
{    }{
}
{   }{
}
}
{}{}{}
ein Untermodul des Restklassenmoduls. Das Urbild von $V$ in $M$ unter der Restklassenabbildung sei $\tilde{V}$. Dieser Modul ist nach Voraussetzung endlich erzeugt, und die Bilder eines solchen Erzeugendensystems erzeugen auch den Bildmodul $V$.

Es seien nun die äußeren Moduln $M_1$ und $M_3$ noethersch, und sei

\mavergleichskette
{\vergleichskette
{ U
}
{ \subseteq }{ M
}
{  }{
}
{    }{
}
{   }{
}
}
{}{}{}
ein Untermodul. Es sei

\mavergleichskette
{\vergleichskette
{ U_3
}
{ \subseteq }{ M_3
}
{  }{
}
{    }{
}
{   }{
}
}
{}{}{}
der Bild-Untermodul davon. $U_3$ wird von endlich vielen Elementen
\mathl{s_1 , \ldots , s_n}{} erzeugt, und wir können annehmen, dass diese

\mavergleichskette
{\vergleichskette
{ s_i
}
{ = }{ \overline{r}_i
}
{  }{
}
{    }{
}
{   }{
}
}
{}{}{}
die Bilder von Elementen

\mavergleichskette
{\vergleichskette
{ r_i
}
{ \in }{ U
}
{  }{
}
{    }{
}
{   }{
}
}
{}{}{}
sind. Betrachte
\mathl{U \cap M_1}{.} Dies ist ein Untermodul von $M_1$, und daher endlich erzeugt, sagen wir von
\mathl{t_1 , \ldots , t_k}{,} die wir als Elemente in $U$ auffassen. Wir behaupten, dass

\mathdisp {r_1 , \ldots , r_n,t_1 , \ldots , t_k} {  }

ein Erzeugendensystem von $U$ bilden. Es sei dazu

\mavergleichskette
{\vergleichskette
{ m
}
{ \in }{ U
}
{  }{
}
{    }{
}
{   }{
}
}
{}{}{}
ein beliebiges Element. Dann ist

\mavergleichskette
{\vergleichskette
{ \overline{m}
}
{ = }{ \sum_{i = 1}^n  a_i s_i
}
{  }{
}
{    }{
}
{   }{
}
}
{}{}{}
und daher geht das Element
\mathl{m-\sum_{i=1}^n  a_i r_i}{} rechts auf $0$. Dann gehört es aber zum Kern der Restklassenabbildung, also zu $M_1$. Andererseits gehört dieses Element auch zu $U$, also zum Durchschnitt
\mathl{M_1 \cap U}{,} der ja von den
\mathl{t_1 , \ldots , t_k}{} erzeugt wird. Also kann man

\mavergleichskettedisp
{\vergleichskette
{ m-\sum_{i = 1}^n  a_i r_i
}
{ =} { \sum_{j = 1}^k  b_j t_j
}
{ } {
}
{ } {
}
{ } {
}
}
{}{}{}
bzw.

\mavergleichskette
{\vergleichskette
{m
}
{ = }{ \sum_{i = 1}^n  a_i r_i+ \sum_{j = 1}^k  b_j t_j
}
{  }{
}
{    }{
}
{   }{
}
}
{}{}{}
schreiben.




}






\inputfaktbeweis
{Kommutative Ringtheorie/Noethersche Ringe/Endlich erzeugte Moduln sind noethersch/Fakt}
{Satz}
{}
{



\faktsituation {}
\faktvoraussetzung {Es sei $R$ ein
\definitionsverweis {noetherscher}{}{}
\definitionsverweis {kommutativer Ring}{}{} und $M$ ein
\definitionsverweis {endlich erzeugter}{}{}
$R$-\definitionsverweis {Modul}{}{.}}
\faktfolgerung {Dann ist $M$ ein
\definitionsverweis {noetherscher Modul}{}{.}}
\faktzusatz {}
\faktzusatz {}





}
{





Wir beweisen die Aussage durch Induktion über die Anzahl $n$ der Modulerzeuger von $M$. Bei

\mavergleichskette
{\vergleichskette
{ n
}
{ = }{ 0
}
{  }{
}
{    }{
}
{   }{
}
}
{}{}{}
liegt der Nullmodul vor. Es sei

\mavergleichskette
{\vergleichskette
{ n
}
{ = }{ 1
}
{  }{
}
{    }{
}
{   }{
}
}
{}{}{.}
Dann gibt es eine surjektive Abbildung
\maabb {} { R } { M \cong R/ {\mathfrak a}
} {.}
Nach

Lemma 10.3

ist aber ein Restklassenmodul eines noetherschen Moduls wieder noethersch, und der Ring selbst ist nach Voraussetzung noethersch, also ist $M$ noethersch.

Es sei nun

\mavergleichskette
{\vergleichskette
{ n
}
{ \geq }{ 2
}
{  }{
}
{    }{
}
{   }{
}
}
{}{}{}
und die Aussage für kleinere $n$ bereits bewiesen. Es sei
\mathl{m_1 , \ldots ,  m_n}{} ein Erzeugendensystem von $M$. Wir betrachten den durch
\mathl{m_1 , \ldots , m_{n-1}}{} erzeugten $R$-Untermodul, den wir mit $M_1$ bezeichnen. Dieser Untermodul gibt Anlass zu einer
\definitionsverweis {kurzen exakten Sequenz}{}{,}
nämlich

\mathdisp {0 \longrightarrow M_1 \longrightarrow M \longrightarrow M/M_1 =:M_3 \longrightarrow 0} { . }

Hier wird der linke Modul von \mathlk{n-1}{} Elementen erzeugt und ist nach Induktionsvoraussetzung noethersch. Der rechte Modul wird von der Restklasse von $m_n$, also von einem Element erzeugt, ist also auch noethersch. Nach

Lemma 10.3

ist dann $M$ noethersch.




}






  \zwischenueberschrift{Hilbertscher Nullstellensatz - algebraische Version}

Im Folgenden werden für eine
$R$-\definitionsverweis {Algebra}{}{}
$A$ die Begriffe
\definitionsverweis {endlich}{}{}
und
\definitionsverweis {endlich erzeugt}{}{}
wichtig sein. Ersteres bedeutet, dass $A$, aufgefasst als
$R$-\definitionsverweis {Modul}{}{,}
endlich erzeugt ist, das zweite, dass $A$ als Algebra endlich erzeugt ist. Der Polynomring
\mathl{R[X_1 , \ldots , X_n]}{} ist im letzteren Sinne endlich erzeugt, und zwar bilden die Variablen ein endliches Algebra-Erzeugendensystem. Der Polynomring ist aber nicht als Modul endlich erzeugt, das einfachste Modul-Erzeugendensystem ist durch alle Monome gegeben.

Wir wollen die algebraische Version des Hilbertschen Nullstellensatzes beweisen. Dazu benötigen wir die beiden folgenden Lemmata.





\inputfaktbeweis
{Endlich erzeugte kommutative Algebren/R noethersch/A über R endlich erzeugt/A endlich über B/B ist endlich erzeugt/Fakt}
{Lemma}
{}
{



\faktsituation {}
\faktvoraussetzung {Es sei $R$ ein
\definitionsverweis {noetherscher}{}{}
\definitionsverweis {kommutativer Ring}{}{}
und $A$ eine
\definitionsverweis {endlich erzeugte}{}{}
$R$-Algebra. Es sei

\mavergleichskette
{\vergleichskette
{ B
}
{ \subseteq }{ A
}
{  }{
}
{    }{
}
{   }{
}
}
{}{}{}
eine $R$-Unteralgebra, über der $A$
\definitionsverweis {endlich}{}{}
\zusatzklammer {als $B$-Modul} {} {}
sei.}
\faktfolgerung {Dann ist auch $B$ eine endlich erzeugte $R$-Algebra.}
\faktzusatz {}
\faktzusatz {}





}
{





Wir schreiben

\mavergleichskette
{\vergleichskette
{A
}
{ = }{ R[x_1 , \ldots ,  x_n]
}
{  }{
}
{    }{
}
{   }{
}
}
{}{}{}
und

\mavergleichskette
{\vergleichskette
{ A
}
{ = }{ Ba_1 + \cdots +  Ba_m
}
{  }{
}
{    }{
}
{   }{
}
}
{}{}{}
mit

\mavergleichskette
{\vergleichskette
{ a_i
}
{ \in }{ A
}
{  }{
}
{    }{
}
{   }{
}
}
{}{}{.}
Wir setzen
\mathkor {} {x_i = \sum_{j=1}^m b_{ij} a_j} {und} {a_ia_j= \sum_{k=1}^m b_{ijk} a_k} {}
mit Koeffizienten

\mavergleichskette
{\vergleichskette
{ b_{ij}, b_{ijk}
}
{ \in }{ B
}
{  }{
}
{    }{
}
{   }{
}
}
{}{}{.}
Wir betrachten die von diesen Koeffizienten erzeugte $R$-Unteralgebra $S$ von $B$ und den
$S$-\definitionsverweis {Untermodul}{}{}

\mavergleichskette
{\vergleichskette
{ \tilde{A}
}
{ = }{ Sa_1 + \cdots + Sa_m
}
{ \subseteq }{ A
}
{    }{
}
{   }{
}
}
{}{}{.}
Die Produkte
\mathl{a_i a_j}{} gehören wieder zu diesem Modul, daher ist $\tilde{A}$ sogar eine $S$-Algebra. Weil die $x_i$ ebenfalls zu $\tilde{A}$ gehören, gilt sogar

\mavergleichskette
{\vergleichskette
{A
}
{ = }{ \tilde{A}
}
{  }{
}
{    }{
}
{   }{
}
}
{}{}{.}
Dies bedeutet, dass $A$ ein endlicher $S$-Modul ist. Nach

Korollar 9.9

ist $S$ ein noetherscher Ring und nach

Satz 10.4

ist der $S$-Untermodul

\mavergleichskette
{\vergleichskette
{B
}
{ \subseteq }{ A
}
{  }{
}
{    }{
}
{   }{
}
}
{}{}{}
ebenfalls endlicher $S$-Modul. Die Kette

\mavergleichskette
{\vergleichskette
{R
}
{ \subseteq }{ S
}
{ \subseteq }{ B
}
{    }{
}
{   }{
}
}
{}{}{}
zeigt schließlich, dass $B$ eine endlich erzeugte $R$-Algebra ist.




}






\inputfaktbeweis
{Endlich erzeugte kommutative Algebren/Rationaler Funktionenkörper ist nicht endlich erzeugt/Fakt}
{Lemma}
{}
{



\faktsituation {}
\faktvoraussetzung {Es sei $K$ ein
\definitionsverweis {Körper}{}{}
und

\mavergleichskette
{\vergleichskette
{ R
}
{ = }{ K(X)
}
{  }{
}
{    }{
}
{   }{
}
}
{}{}{}
der zugehörige
\definitionsverweis {rationale Funktionenkörper}{}{.}}
\faktfolgerung {Dann ist $R$ keine endlich erzeugte $K$-Algebra.}
\faktzusatz {}
\faktzusatz {}





}
{





Es sei angenommen, dass die rationalen Funktionen

\mathbed {F_i= \frac{P_i }{Q_i }} {}
 {i=1, \ldots , n} {}
 {} {} {} {,}
ein endliches Erzeugendensystem von
\mathl{K(X)}{} bilden, mit

\mathbed {P_i,Q_i \in K[X]} {}
 {Q_i \neq 0} {}
 {} {} {} {.}
Durch Übergang zu einem Hauptnenner kann man annehmen, dass die Nenner

\mavergleichskette
{\vergleichskette
{Q_i
}
{ = }{ Q
}
{  }{
}
{    }{
}
{   }{
}
}
{}{}{}
gleich sind. Die Annahme bedeutet also insbesondere, dass der Körper der rationalen Funktionen sich durch Nenneraufnahme an nur einem Element ergeben würde. Da $Q$ keine Konstante ist
\zusatzklammer {sonst wäre

\mavergleichskettek
{\vergleichskettek
{ K[X]
}
{ = }{ K(X)
}
{  }{
}
{    }{
}
{   }{
}
}
{}{}{,}
was nicht der Fall ist} {} {,}
ist

\mavergleichskette
{\vergleichskette
{Q-1
}
{ \neq }{ 0
}
{  }{
}
{    }{
}
{   }{
}
}
{}{}{}
und daher ist

\mavergleichskette
{\vergleichskette
{ \frac{1}{Q-1}
}
{ \in }{ K(X)
}
{  }{
}
{    }{
}
{   }{
}
}
{}{}{.}
Also gibt es eine Darstellung

\mavergleichskettedisp
{\vergleichskette
{ \frac{1}{Q-1}
}
{ =} { \frac{P}{Q^s}
}
{ } {
}
{ } {
}
{ } {
}
}
{}{}{}
mit einem geeigneten $s$. Daraus folgt

\mavergleichskette
{\vergleichskette
{Q^s
}
{ = }{ (Q-1)P
}
{  }{
}
{    }{
}
{   }{
}
}
{}{}{.}
Da
\mathkor {} {Q^s} {und} {Q-1} {}
das
\definitionsverweis {Einheitsideal}{}{}
in
\mathl{K[X]}{} erzeugen, folgt daraus, dass bereits
\mathl{Q-1}{} das Einheitsideal erzeugt, also selbst eine Einheit ist. Dann wäre $Q$ aber doch eine Konstante, was es nicht ist.




}



Die folgende Aussage ist die algebraische Version des Hilbertschen Nullstellensatzes.



\inputfaktbeweis
{Hilbertscher Nullstellensatz (algebraisch)/Endlich erzeugte Körpererweiterung ist endlich/Fakt}
{Satz}
{}
{



\faktsituation {}
\faktvoraussetzung {Es sei $K$ ein
\definitionsverweis {Körper}{}{}
und sei

\mavergleichskette
{\vergleichskette
{K
}
{ \subseteq }{ L
}
{  }{
}
{    }{
}
{   }{
}
}
{}{}{}
eine
\definitionsverweis {Körpererweiterung}{}{,}
die
\zusatzklammer {als $K$-Algebra} {} {}
\definitionsverweis {endlich erzeugt}{}{}
sei.}
\faktfolgerung {Dann ist $L$
\definitionsverweis {endlich}{}{}
über $K$.}
\faktzusatz {}
\faktzusatz {}





}
{





Wir setzen

\mavergleichskette
{\vergleichskette
{ L
}
{ = }{ K[x_1 , \ldots ,  x_n]
}
{  }{
}
{    }{
}
{   }{
}
}
{}{}{.}
Es sei $K_i$ der
\definitionsverweis {Quotientenkörper}{}{}
von
\mathl{K[x_1 , \ldots , x_i]}{}
\zusatzklammer {innerhalb von $L$} {} {.}
Wir haben also eine Körperkette

\mavergleichskettedisp
{\vergleichskette
{ K
}
{ =} { K_0
}
{ \subseteq} { K_1
}
{ \subseteq} { \ldots
}
{ \subseteq} { K_n
}
}
{
\vergleichskettefortsetzung
{ =} { L
}
{ } {}
{ } {}
{ } {}
}{}{.}
Wir wollen zeigen, dass $L$ endlich über $K$ ist, und dazu genügt es
nach
Satz 2.8 (Körper- und Galoistheorie (Osnabrück 2018-2019))

zu zeigen, dass jeder Schritt in der Körperkette endlich ist. Es sei angenommen, dass

\mavergleichskette
{\vergleichskette
{ K_i
}
{ \subseteq }{ K_{i+1}
}
{  }{
}
{    }{
}
{   }{
}
}
{}{}{}
nicht endlich ist, aber alle folgenden Schritte endlich sind. Wir wenden

Lemma 10.5

auf

\mavergleichskettedisp
{\vergleichskette
{K
}
{ \subseteq} { K_{i+1}
}
{ \subset} { L
}
{ } {
}
{ } {
}
}
{}{}{}
an und erhalten, dass
\mathl{K_{i+1}}{}
\definitionsverweis {endlich erzeugt}{}{}
über $K$ ist. Dann ist insbesondere
\mathl{K_{i+1}}{} auch endlich erzeugt über $K_i$. Andererseits ist
\mathl{K_{i+1}}{} der Quotientenkörper von
\mathl{K_i[x_{i+1}]}{.} Wir haben also eine Kette

\mavergleichskettedisp
{\vergleichskette
{ K_i
}
{ \subseteq} { K_i[x_{i+1}]
}
{ \subseteq} { Q( K_i[x_{i+1}])
}
{ =} { K_{i+1}
}
{ } {
}
}
{}{}{,}
wo
\mathl{K_{i+1}}{} endlich erzeugt über $K_i$ ist, aber nicht endlich. Wäre
\mathl{x_{i+1}}{}
\definitionsverweis {algebraisch}{}{}
über $K_i$, so auch endlich, und dann wäre
\mathl{K_i[x_{i+1}]}{} bereits ein Körper
nach
Aufgabe 10.1
.
Dann wäre die letzte Kette insgesamt endlich, im Widerspruch zur Wahl von $i$. Also ist
\mathl{x_{i+1}}{}
\definitionsverweis {transzendent}{}{}
über $K_i$. Dann ist aber
\mathl{K_i[x_{i+1}]}{} isomorph zu einem
\definitionsverweis {Polynomring in einer Variablen}{}{}
und
\mathl{Q( K_i[x_{i+1}])}{} ist isomorph zum
\definitionsverweis {rationalen Funktionenkörper}{}{}
über $K_i$. Dieser ist aber nach

Lemma 10.6

nicht endlich erzeugt, sodass sich erneut ein Widerspruch ergibt.




}






\inputfaktbeweis
{Algebren von endlichem Typ über Körper/Homomorphismen/Urbild von maximalem Ideal ist maximal/Fakt}
{Satz}
{}
{



\faktsituation {}
\faktvoraussetzung {Es sei $K$ ein
\definitionsverweis {Körper}{}{} und seien $A$ und $B$ zwei $K$-Algebren von
\definitionsverweis {endlichem Typ}{}{.}
Es sei
\maabb {\varphi} { A } { B
} {}
ein
$K$-\definitionsverweis {Algebrahomomorphismus}{}{.}}
\faktfolgerung {Dann ist für jedes
\definitionsverweis {maximale Ideal}{}{}
${\mathfrak m}$ aus $B$ auch das Urbild
\mathl{\varphi ^{-1} ({\mathfrak m})}{} ein maximales Ideal.}
\faktzusatz {}
\faktzusatz {}





}
{





Es sei ${\mathfrak m}$ ein maximales Ideal aus $B$. Wir wissen nach

Aufgabe 4.19
,
dass unter jedem
\definitionsverweis {Ringhomomorphismus}{}{}
das Urbild eines
\definitionsverweis {Primideals}{}{}
wieder prim ist, also ist
\mathl{\varphi^{-1}( {\mathfrak m} )}{} zunächst ein Primideal, das wir ${\mathfrak p}$ nennen. Wir erhalten induzierte Ringhomomorphismen

\mathdisp {K \longrightarrow A/ {\mathfrak p} \longrightarrow B/ {\mathfrak m} = L} { , }

wobei $L$ ein Körper ist und wobei beide Homomorphismen injektiv und von endlichem Typ sind. Da die Gesamtabbildung von endlichem Typ ist und $K$ und $L$ Körper sind, folgt aus

Satz 10.7
,
dass diese Abbildung
\definitionsverweis {endlich}{}{}
ist. Wir wollen zeigen, dass der Zwischenring
\mathl{A/ {\mathfrak p}}{} ein
\definitionsverweis {Körper}{}{}
ist. Dies folgt aber aus

Aufgabe 10.2
.




}






\inputfaktbeweis
{Algebren von endlichem Typ über Körper/Radikal ist Durchschnitt von maximalen Idealen/Fakt}
{Satz}
{}
{



\faktsituation {}
\faktvoraussetzung {Es sei $K$ ein
\definitionsverweis {Körper}{}{} und sei $A$ eine $K$-Algebra von
\definitionsverweis {endlichem Typ}{}{.}}
\faktfolgerung {Dann ist jedes
\definitionsverweis {Radikal}{}{} in $A$ der Durchschnitt von
\definitionsverweis {maximalen Idealen}{}{.}}
\faktzusatz {}
\faktzusatz {}





}
{





Nach
Aufgabe 10.17

ist jedes Radikal der Durchschnitt von Primidealen. Es genügt also zu zeigen, dass jedes Primideal in einer endlich erzeugten Algebra der Durchschnitt von maximalen Idealen ist. Es sei ${\mathfrak p}$ ein Primideal und

\mavergleichskette
{\vergleichskette
{ f
}
{ \notin }{ {\mathfrak p}
}
{  }{
}
{    }{
}
{   }{
}
}
{}{}{.}
Dann ist ${\mathfrak p}$ ein Primideal in der Nenneraufnahme

\mavergleichskette
{\vergleichskette
{ B
}
{ \defeq }{ A_f
}
{  }{
}
{    }{
}
{   }{
}
}
{}{}{.}
Es gibt ein
\zusatzklammer {in $A_f$} {} {}
maximales Ideal

\mavergleichskette
{\vergleichskette
{ {\mathfrak m}
}
{ \subset }{ A_f
}
{  }{
}
{    }{
}
{   }{
}
}
{}{}{}
oberhalb von
\mathl{{\mathfrak p} A_f}{.} Wir fassen $A_f$ als endlich erzeugte $K$-Algebra auf und betrachten
\maabbdisp {\varphi} { A } { A_f
} {.}
Dann ist

\mavergleichskette
{\vergleichskette
{ {\mathfrak p}
}
{ \subseteq }{ \varphi^{-1} ({\mathfrak m})
}
{  }{
}
{    }{
}
{   }{
}
}
{}{}{}
und

\mavergleichskette
{\vergleichskette
{ f
}
{ \notin }{ \varphi^{-1} ({\mathfrak m})
}
{  }{
}
{    }{
}
{   }{
}
}
{}{}{.}
Nach

Satz 10.8

ist
\mathl{\varphi^{-1} ({\mathfrak m})}{} maximal.




}






\inputfaktbeweis
{Algebren von endlichem Typ über Körper/Algebraisch abgeschlossen/Maximale Ideale sind Punktideal/Fakt}
{Satz}
{}
{



\faktsituation {}
\faktvoraussetzung {Es sei $K$ ein
\definitionsverweis {algebraisch abgeschlossener Körper}{}{}
und sei $A$ eine
\definitionsverweis {endlich erzeugte}{}{}
$K$-\definitionsverweis {Algebra}{}{.}}
\faktfolgerung {Dann ist jeder
\definitionsverweis {Restklassenkörper}{}{}
von $A$
\definitionsverweis {isomorph}{}{}
zu $K$.}
\faktzusatz {Anders formuliert: Jedes
\definitionsverweis {maximale Ideal}{}{}
in $A$ ist ein
\definitionsverweis {Punktideal}{}{.}}
\faktzusatz {}





}
{





Es sei ${\mathfrak m}$ ein maximales Ideal der endlich erzeugten $K$-Algebra $A$ und betrachte

\mathdisp {K \longrightarrow A \longrightarrow A/{\mathfrak m} \defeqr  L} { . }

Hier ist $L$ ein Körper und zugleich eine endlich erzeugte $K$-Algebra. Nach

Satz 10.7

muss also $L$ eine endliche $K$-Algebra sein. Da $K$ algebraisch abgeschlossen ist, muss

\mavergleichskette
{\vergleichskette
{ K
}
{ = }{ L
}
{  }{
}
{    }{
}
{   }{
}
}
{}{}{}
sein.




}
